\chapter{Male Reproductive System}

\section*{Patient focus}
This chapter explains the examination or clinical problem in accessible language. Patients should use it to prepare questions and understand the purpose, experience, benefits, and limitations of imaging; local instructions and the official report take priority.


\label{ch:ch51_male_reproductive}

\section{What imaging covers}
The male reproductive system includes the prostate gland, seminal vesicles, testes, epididymis, vas deferens, and penis. Imaging is most commonly requested for:
\begin{itemize}
    \item \textbf{Testicular symptoms:} Acute pain (torsion, epididymitis), a lump, or swelling.
    \item \textbf{Prostate cancer:} Staging, localisation for biopsy, and treatment planning.
    \item \textbf{Infertility evaluation:} Assessment of seminal vesicles, vas deferens, and varicocele.
    \item \textbf{Penile trauma or suspected Peyronie disease.}
\end{itemize}

\section{What happens during the scan?}

\subsection{Testicular ultrasound}
This is the first-line investigation for acute scrotal symptoms.
\begin{itemize}
    \item A high-frequency linear probe is placed directly on the scrotum with warm gel.
    \item Both testes are compared side by side.
    \item Colour and spectral Doppler assess blood flow.
    \item The scan takes 10--20 minutes.
    \item No preparation is required.
\end{itemize}

\subsection{Prostate MRI (multiparametric)}
This is the standard for prostate cancer localisation and staging.
\begin{itemize}
    \item Performed on a 1.5T or 3T scanner.
    \item Involves T2-weighted imaging, diffusion-weighted imaging (DWI), and dynamic contrast-enhanced (DCE) sequences.
    \item An endorectal coil may be used at 1.5T but is not required at 3T.
    \item Takes 30--45 minutes.
    \item Rectal gel or an antispasmodic may be given to reduce motion artefact.
\end{itemize}

\section{Testicular ultrasound: normal and abnormal}

\subsection{Normal appearance}
The normal testis has homogeneous medium-level echogenicity, surrounded by the tunica albuginea (a thin echogenic line). The epididymis is seen as a comma-shaped structure at the posterior pole. Blood flow is symmetric between both testes.

\subsection{Torsion}
Testicular torsion is a surgical emergency where the spermatic cord twists, cutting off blood supply.
\begin{itemize}
    \item \textbf{Early ($<$6 hours):} Testis may appear normal in echogenicity; Doppler shows absent or markedly reduced flow.
    \item \textbf{Intermediate (6--12 hours):} Testis becomes hypoechoic; no flow on colour Doppler.
    \item \textbf{Late ($>$12 hours):} Hypoechoic or heterogeneous testis; may be enlarged and swollen.
    \item \textbf{Incomplete torsion:} Colour Doppler shows asymmetrically reduced flow---this does \emph{not} exclude torsion.
\end{itemize}

\begin{pearlbox}[Don't wait for the ultrasound]
In a male with acute scrotal pain and clinical signs of torsion (absent cremasteric reflex, high-riding testis, abnormal lie), surgical exploration should not be delayed for imaging. If doubt exists, proceed to ultrasound---but if it is equivocal, proceed to the operating theatre. Time is testis.
\end{pearlbox}

\subsection{Epididymitis and orchitis}
Epididymitis is the most common cause of acute scrotal pain in adult males. Ultrasound shows:
\begin{itemize}
    \item Enlarged, hypoechoic epididymis
    \item Increased blood flow on Doppler
    \item Associated reactive hydrocele
    \item If the testis is also involved (orchitis), it becomes hypoechoic with increased vascularity
\end{itemize}

\subsection{Testicular masses}
Solid intratesticular masses are testicular cancer until proven otherwise in adults. Ultrasound features:
\begin{itemize}
    \item Most testicular cancers are \textbf{germ cell tumours} (seminoma or non-seminomatous).
    \item Seminomas tend to be homogeneous, hypoechoic, and well-defined.
    \item Non-seminomatous tumours are often heterogeneous, with cystic or necrotic areas.
    \item \textbf{Lymphoma:} Typically in men $>$60 years; bilateral in 10--25\%; may be diffusely infiltrative.
    \item \textbf{Metastases:} Lymphoma, leukaemia, prostate, lung.
    \item All solid intratesticular masses require urological referral and tumour markers ($\alpha$-FP, $\beta$-hCG, LDH).
\end{itemize}

\section{Varicocele}
A varicocele is dilatation of the pampiniform venous plexus, seen as a cluster of serpentine, anechoic tubular structures above and behind the testis.
\begin{itemize}
    \item Occurs in 15\% of males, more common on the left (due to left gonadal vein drainage into the left renal vein).
    \item Visible on ultrasound $>$3~mm in diameter.
    \item Doppler shows venous flow that increases with Valsalva manoeuvre.
    \item Clinically significant when associated with infertility, pain, or testicular atrophy.
\end{itemize}

\section{Prostate MRI: PI-RADS v2.1 scoring}
The Prostate Imaging Reporting and Data System (PI-RADS) provides a standardised language for reporting prostate MRI:

\begin{table}[h]
\centering
\caption{PI-RADS v2.1 scoring system}
\label{tab:pirads}
\begin{tabular}{clll}
\hline
\textbf{Score} & \textbf{Term} & \textbf{Meaning} & \textbf{Clinical action} \\
\hline
1 & Very low & Clinically significant cancer extremely unlikely & No biopsy \\
2 & Low & Clinically significant cancer unlikely & Consider biopsy based on risk \\
3 & Intermediate & Equivocal; indeterminate risk & Biopsy may be recommended \\
4 & High & Clinically significant cancer likely & MRI-targeted biopsy recommended \\
5 & Very high & Clinically significant cancer extremely likely & MRI-targeted biopsy recommended \\
\hline
\end{tabular}
\end{table}

\section{PI-RADS assessment: dominant sequence}
The scoring depends on lesion location (peripheral zone vs transition zone):

\begin{table}[h]
\centering
\caption{PI-RADS dominant sequence by zone}
\label{tab:pirads_zone}
\begin{tabular}{ll}
\hline
\textbf{Location} & \textbf{Dominant sequence} \\
\hline
Both zones & DCE is used as a tiebreaker (positive = upgrade by 1 point) \\
Peripheral zone & DWI (apparent diffusion coefficient mapping) \\
Transition zone & T2-weighted imaging \\
\hline
\end{tabular}
\end{table}

\begin{pearlbox}[When to recommend MRI vs biopsy]
Multiparametric MRI is recommended \emph{before} first prostate biopsy in men with elevated PSA (NICE, EAU guidelines). A PI-RADS $\geq$4 lesion warrants MRI-targeted biopsy (in-bore or cognitive fusion). PI-RADS 3 may warrant biopsy depending on clinical risk (PSA density, family history, ethnicity). PI-RADS 1--2 can usually reassure against immediate biopsy.
\end{pearlbox}

\section{Prostate MRI: staging findings}
For known prostate cancer, MRI staging evaluates:
\begin{itemize}
    \item \textbf{T2 (organ-confined):} Low T2 signal within the capsule, no extraprostatic extension.
    \item \textbf{T3a (extraprostatic extension):} Ill-defined tumour extending beyond the capsule, neurovascular bundle involvement.
    \item \textbf{T3b (seminal vesicle invasion):} Loss of normal high T2 signal in the seminal vesicles with enhancing tumour.
    \item \textbf{T4 (invasion of adjacent structures):} Bladder, rectum, or pelvic sidewall involvement.
    \item \textbf{Lymph nodes:} Short-axis diameter $>$1~cm or with loss of normal fatty hilum.
\end{itemize}

\section{Scrotal masses: extratesticular vs intratesticular}
Location is the single most important distinction:

\begin{table}[h]
\centering
\caption{Extratesticular vs intratesticular scrotal masses}
\label{tab:scrotal_masses}
\begin{tabular}{lll}
\hline
\textbf{Intratesticular} & \textbf{Extratesticular} & \textbf{Key point} \\
\hline
Germ cell tumour & Spermatocele & Extratesticular masses are usually benign \\
Leukaemia & Lipoma & Consider haematological malignancy \\
Lymphoma & Adenomatoid tumour & Intratesticular solid masses are malignant until proven otherwise \\
\hline
\end{tabular}
\end{table}

\section{Penile imaging}
Penile imaging is uncommon but indicated in:
\begin{itemize}
    \item \textbf{Priapism:} Colour Doppler to assess arterial flow (ischaemic vs high-flow).
    \item \textbf{Peyronie disease:} Ultrasound with Doppler to identify the fibrous plaque and assess vascularity.
    \item \textbf{Penile fracture:} MRI to demonstrate tunica albuginea disruption and urethral injury.
    \item \textbf{Penile cancer:} MRI for local staging (depth of invasion, involvement of corpora, urethra).
\end{itemize}

\section{Reporting checklist}
A structured report for male reproductive imaging should include:
\begin{enumerate}
    \item Testes: size, echogenicity, symmetry, focal lesions (intratesticular vs extratesticular)
    \item Blood flow: symmetric, asymmetric, absent
    \item Epididymis: size, echogenicity, vascularity
    \item Scrotal wall, hydrocele, haematocele
    \item Prostate: size, zonal anatomy, PI-RADS score (if MRI)
    \item Seminal vesicles: symmetry, signal, cysts
    \item Lymph nodes
    \item Comparison with prior imaging
\end{enumerate}

\clearpage
